\section{The game, and the rule dialect studied}
\label{sec:rules}

Canadian Fish (also called Literature) is played by six players in two teams of
three, seated so that the teams alternate. The deck is the standard 52 cards
plus two distinguishable jokers, partitioned into nine \emph{half-suits} of six
cards: the low band $\{2,3,4,5,6,7\}$ and the high band $\{9,10,J,Q,K,A\}$ of
each of the four suits, together with a ninth half-suit containing the four
eights and the two jokers. Every player is dealt nine cards. The team that
claims at least five of the nine half-suits wins.

\paragraph{Asking.} On their turn a player asks one opponent for one named card.
The ask is legal only if the asker holds at least one \emph{other} card of that
card's half-suit, does not hold the named card, and the target still has cards.
If the target holds the card they must surrender it and the asker continues; if
not, the turn passes to the target. Every ask, its answer, and every resulting
transfer is public, as is every player's card count.

\paragraph{Declaring.} A player may declare a half-suit \emph{at any moment},
including during an opponent's turn, and needs to hold none of its cards to do
so. The declaration must name, for each of the six cards, the exact teammate who
holds it. If every card and every assignment is right, the declaring team scores
the half-suit; \emph{any} error awards it to the opposing team. Either way the
six cards leave the game, and the turn returns to whoever held it.

\paragraph{Running out.} A player with no cards drops out of the asking cycle:
they can neither ask nor be asked. If the turn reaches them --- which can only
happen after a declaration --- they hand it to a teammate who still has cards,
and the team may negotiate openly, but may communicate nothing beyond
willingness to receive. If an entire team is cardless, the other team must
declare every remaining half-suit with no further asking, again choosing the
declarer by open negotiation limited to willingness.

\subsection{What the v0.3 simulator modelled, and what it did not}

The v0.3 engine \cite{fishbot03} implemented the deck, the ask legality
conditions, transfers, exact-allocation declarations and the cardless-team
endgame, and it is the direct ancestor of this work. Comparing it against the
authoritative rules above surfaces four gaps, all of which change strategy
rather than merely presentation:

\begin{enumerate}
  \item \textbf{Declarations were restricted to the declarer's own turn.} The
        real rule permits declaring at any moment. This matters because it turns
        declaration timing into a continuous decision rather than a once-per-turn
        one, and because it allows a team to cash a half-suit immediately before
        an opponent could act.
  \item \textbf{The cardless player's choice of successor was fixed} to the
        teammate holding the most cards, rather than being a decision.
  \item \textbf{The forced endgame selected the declarer by comparing private
        confidences}, which leaks more than the rules permit; the rules allow the
        team to exchange only willingness.
  \item \textbf{Games were adjudicated by majority holding} after an action cap.
        A cap is a necessary safety device --- \S\ref{sec:termination} shows the
        rules do not guarantee termination against deterministic play --- but
        adjudication changes the payoff structure of exactly the stalling
        behaviour that strong policies produce. v0.4 makes termination a policy
        property instead, and retains a cap only as a backstop whose residue is
        split neutrally by majority holding and whose incidence is reported with
        every experiment.
\end{enumerate}

The v0.4 engine implements all four correctly and exposes each as a switch, so
that the effect of the rule dialect is measurable rather than assumed
(\S\ref{sec:robustness}). Unless stated otherwise, every number in this paper
uses the full rules: out-of-turn declarations enabled, cardless players allowed
to declare, willingness-only endgame negotiation, and a $400$-ask safety cap whose incidence we report
($360$ under the v0.3 dialect, matching that engine).
